\documentclass[a4paper,11pt]{article}
\usepackage[T2A]{fontenc}
\usepackage[utf8]{inputenc}
\usepackage[russian]{babel}
\usepackage{amsmath,amssymb,mathtools}
\usepackage{cmbright}
\renewcommand{\familydefault}{\sfdefault}
\usepackage{icomma}
\usepackage[a4paper,left=16mm,right=16mm,top=12mm,bottom=11mm]{geometry}
\usepackage{microtype}
\usepackage{multicol}
\usepackage{xcolor}
\usepackage{hyperref}
\definecolor{accent}{HTML}{15747A}
\definecolor{darkaccent}{HTML}{164C55}
\definecolor{textgray}{HTML}{283238}
\definecolor{muted}{HTML}{65747A}
\definecolor{rulegray}{HTML}{CADADC}
\color{textgray}
\hypersetup{hidelinks}
\setlength{\parindent}{0pt}
\setlength{\parskip}{0pt}
\setlength{\columnsep}{8mm}
\setlength{\abovedisplayskip}{1mm}
\setlength{\belowdisplayskip}{1mm}
\emergencystretch=2em
\newcommand{\sect}[1]{\par\vspace{2.2mm}{\normalsize\bfseries\color{darkaccent}#1}\par\vspace{0.3mm}{\color{rulegray}\rule{\linewidth}{0.45pt}}\par\vspace{0.8mm}}
\newcommand{\fml}[1]{\par\noindent{\normalsize\(\displaystyle #1\)}\par\vspace{2.8mm}}
\begin{document}
\thispagestyle{empty}

\hypersetup{pdftitle={Главные формулы. София Кальней},pdfauthor={Лёвин Артём Александрович}}
\begin{center}
{\small\bfseries\color{accent}УЧЕБНАЯ ПАМЯТКА \quad\textbullet\quad ОСНОВНЫЕ ФОРМУЛЫ}\par
\vspace{1mm}
{\fontsize{21}{25}\selectfont\bfseries\color{darkaccent} Планиметрия: главные формулы\par}
\vspace{0.8mm}
{\small\color{muted}София Кальней\quad\textbullet\quad ОГЭ · математика\quad\textbullet\quad 12.10.2026}
\end{center}
\vspace{2mm}
\noindent{\color{accent}\rule{\linewidth}{0.9pt}}
\vspace{1.3mm}
{\bfseries\color{darkaccent}Главные принципы решения}\par
\vspace{1mm}
\begin{multicols}{2}
\raggedcolumns
\small
\noindent\textbullet\quad В геометрии сначала ищите прямоугольный треугольник и подобие.\par\vspace{0.75mm}
\noindent\textbullet\quad Для площади выбирайте известные сторону и высоту либо две стороны и угол.\par\vspace{0.75mm}
\columnbreak
\noindent\textbullet\quad Равные касательные из одной точки дают равные отрезки.\par\vspace{0.75mm}
\noindent\textbullet\quad В сложной задаче делайте дополнительное построение, раскрывающее отношение или площадь.\par\vspace{0.75mm}
\end{multicols}
\vspace{1.1mm}
\noindent{\color{accent}\rule{\linewidth}{0.5pt}}
\vspace{0.6mm}
\begin{multicols}{2}
\raggedcolumns
\normalsize
\sect{Треугольник}
\fml{a^2+b^2=c^2\quad(c\ \text{гипотенуза})}
\fml{S=\frac12 ah=\frac12 ab\sin\gamma}
\fml{p=\frac{a+b+c}{2}}
\fml{S=\sqrt{p(p-a)(p-b)(p-c)}}
\sect{Подобие и биссектриса}
\fml{\frac{AB}{A_1B_1}=\frac{BC}{B_1C_1}=k}
\fml{\frac{S_1}{S_2}=k^2}
\fml{\frac{AD}{DC}=\frac{AB}{BC}\;(BD\ \text{бисс.})}
\columnbreak
\sect{Трапеция и окружность}
\fml{S_{\text{трап}}=\frac{a+b}{2}\,h}
\fml{PA=PB\quad\text{(касательные)}}
\fml{\angle AOB=2\angle ACB}
\sect{Прямоугольный треугольник}
\fml{\sin\alpha=\frac{a}{c}}
\fml{\cos\alpha=\frac{b}{c}}
\fml{\tg\alpha=\frac ab}
\end{multicols}
\vfill
{\color{accent}\rule{\linewidth}{0.8pt}}\par\vspace{1.5mm}
{\small\color{muted}Лёвин Артём Александрович эксклюзивно для Софии Кальней}
\end{document}